\hmsection{Conclusion}{}{}

This project developed an autonomous robotic system for
detecting, classifying, and sorting coloured balls using
computer vision and a 4-DOF robotic arm. The system integrates
a deterministic HSV-based detection pipeline running on a
Raspberry Pi~5, a closed-form geometric inverse kinematics
solver, an adaptive current-sensing grip verification mechanism,
and a dedicated firmware layer on the OpenRB-150 microcontroller.

The minimum viable product defined in Section~3.3 was
\textbf{[achieved / partially achieved]}: the system completed
\textbf{[TBD]} consecutive autonomous sorting cycles with a pick
success rate of \textbf{[TBD]}\% and a placement accuracy of
\textbf{[TBD]}\% under standard laboratory lighting conditions
(Section~10.5). \textbf{[If partially achieved, state which
requirements were met and which were not.]}

Several design decisions proved effective. The closed-form
geometric IK solver provided deterministic, sub-millisecond
solutions without the convergence issues of numerical methods.
The adaptive grip algorithm, which exploits the Dynamixel
motor's current and load registers as an impedance-based force
sensor, reliably distinguished between a gripped ball and an
empty close without requiring external sensor hardware. The
touch calibration procedure superseded the polynomial sag model
by providing per-point surface heights that account for the
arm's real-world droop at each calibrated position.

The principal limitations identified during the project are the
elbow motor's thermal load under continuous operation at the scan
pose, the absence of mid-approach visual correction due to the
wrist-mounted camera design, and the four-degree-of-freedom
constraint that restricts the gripper to a fixed downward
orientation. The most impactful improvement identified is to
reduce the mass of the wrist-to-claw assembly, which would
directly lower the gravitational torque on the elbow motor and
reduce the arm's overall sag (Section~11.3).
